\documentclass{article}
\pdfpagewidth=8.5in
\pdfpageheight=11in

\usepackage{ijcai26}
\usepackage{times}
\usepackage{url}
\usepackage[hidelinks]{hyperref}
\usepackage[utf8]{inputenc}
\usepackage{graphicx}
\usepackage{amsmath}
\usepackage{booktabs}
\usepackage[brazilian]{babel}

\urlstyle{same}

% Portuguese abstract heading. The IJCAI style writes "Abstract" directly, so we
% redefine it. Babel already translates "References".
\renewenvironment{abstract}{\centerline{\Large\bf
Resumo}\vspace{0.5ex}\begin{quote}}{\par\end{quote}\vskip 1ex}

% Required by the IJCAI template. Leave this block untouched.
\pdfinfo{
/TemplateVersion (IJCAI.2026.0)
}

% TODO: troque o título e o nível entregue (2 ou 3).
\title{Resolvendo o ARC-AGI-1 com LLMs: Fim a Fim e Geração de Código\\(Nível 2)}

% TODO: o seu nome e o seu e-mail.
\author{
    Nome Sobrenome
    \affiliations
    Especialização em Inteligência Artificial Avançada, UFRGS
    \emails
    seu.email@exemplo.com
}

\begin{document}

\maketitle

\begin{abstract}
    % TODO: de três a cinco frases: problema, o que foi feito, principal resultado e conclusão.
    Escreva aqui o resumo do trabalho.
\end{abstract}

\section{Introdução}

% TODO: o que é o ARC-AGI-1, por que ele é difícil para LLMs e o que este trabalho compara.
O ARC-AGI-1 avalia a capacidade de inferir regras abstratas a partir de poucos
exemplos~\cite{chollet2019arXiv}. Neste trabalho, usamos o Gemini~\cite{geminiteam2025arXiv} \ldots

\section{Método}

% TODO: como o solver funciona em cada abordagem (fim a fim e geração de código), quais
% prompts foram usados e, no nível 3, como funciona o laço de correção.

\section{Experimentos}

% TODO: modelo, nível de raciocínio, as 20 tarefas e o agente de codificação usado.

\begin{table}[h]
  \centering
  \small
  \setlength{\tabcolsep}{3pt}
  \begin{tabular}{lrrrrr}
    \toprule
    Abordagem   & Resolvidas & Entrada & Racioc. & Saída & Tempo (s) \\
    \midrule
    Fim a fim   & --/20      & --      & --         & --    & --        \\
    Código      & --/20      & --      & --         & --    & --        \\
    \bottomrule
  \end{tabular}
  \caption{Tarefas resolvidas, total de tokens e tempo por abordagem.}
  \label{tab:resultados}
\end{table}

\section{Resultados}

% TODO: o que a Tabela~\ref{tab:resultados} mostra, um exemplo de acerto e um de erro.

\section{Conclusão}

% TODO: o que funcionou melhor, a que custo, e uma frase sobre como a IA foi usada.

%% The file named.bst is a bibliography style file for BibTeX 0.99c
\bibliographystyle{named}
\bibliography{referencias}

\end{document}
